\documentclass[11pt,a4paper]{article}
\usepackage[utf8]{inputenc}
\usepackage{amsmath,amssymb}
\usepackage{geometry}
\geometry{margin=2.5cm}
\usepackage{array}

\newcommand{\ud}{\mathrm{d}}
\newcommand{\ssc}[1]{\texttt{#1}}

\title{Modelli Simscape per la dinamica di commutazione:\\
\ssc{diode\_rr} e \ssc{igbt\_dyn}}
\author{}
\date{Luglio 2026 --- rev.\ 1.1}

\begin{document}
\maketitle

\section{Introduzione e scopo}

I due componenti descritti in questo documento servono allo studio dei
transitori di accensione e spegnimento di una cella di commutazione a
IGBT: sovratensione allo spegnimento, reverse recovery del diodo di
ricircolo, ringing con le induttanze parassite di loop, e quindi il
dimensionamento di eventuali snubber. Non calcolano perdite e non hanno
dipendenza dalla temperatura: tutti i parametri vanno estratti dalle
curve del datasheet alle condizioni di interesse (tipicamente le curve a
caldo, $125\,^{\circ}\mathrm{C}$ o $150\,^{\circ}\mathrm{C}$).

Il diodo di ricircolo \`e un componente separato (\ssc{diode\_rr}) da
affiancare esternamente all'IGBT (\ssc{igbt\_dyn}); questo permette di
usare gli stessi blocchi in configurazione chopper, half-bridge o double
pulse test senza duplicare il modello.

Entrambi i componenti sono scritti in linguaggio Simscape in forma
completamente continua ($C^1$), senza \emph{if/else} n\'e eventi: le non
linearit\`a a tratti sono sostituite da raccordi smooth
(sez.~\ref{sec:smooth}). Questo riduce gli zero crossing e rende il
modello adatto ai solver a passo variabile.

\section{Funzioni di regolarizzazione numerica}
\label{sec:smooth}

Tutte le funzioni a tratti sono costruite a partire dalla \emph{parte
positiva regolarizzata}
\begin{equation}
p_\delta(x) \;=\; \tfrac{1}{2}\left(x + \sqrt{x^2 + \delta^2}\right)
\;\approx\; \max(x,\,0),
\label{eq:softplus}
\end{equation}
che \`e $C^\infty$, coincide con $\max(x,0)$ per $|x|\gg\delta$ e vale
$\delta/2$ in $x=0$. Da questa si ottengono
\begin{align}
\min\nolimits_\delta(a,b) &= \tfrac{1}{2}\left(a+b-\sqrt{(a-b)^2+\delta^2}\right),\\
\min\nolimits_\delta(x,\,x_{\max}) &= x_{\max} - p_\delta(x_{\max}-x).
\label{eq:smoothmin}
\end{align}
Il parametro $\delta$ (nei componenti: \ssc{dsm}) fissa la larghezza
della zona di raccordo; il default \`e $10\,$mV per il diodo e
$0{,}1\,$V per l'IGBT, trascurabili rispetto alle escursioni di tensione
in gioco. Se il solver fatica, allargare \ssc{dsm} \`e il primo
intervento da provare, prima di toccare le tolleranze.

\paragraph{Esponenziale limitato.}
Nelle iterazioni di Newton il solver valuta le equazioni anche in punti
non fisici (ad esempio $v_j$ di decine di volt in diretta); l'esponenziale
della legge di giunzione va in overflow, il residuo diventa \emph{NaN} e
il passo fallisce. Per questo l'argomento dell'esponenziale \`e limitato
in forma smooth,
\begin{equation}
x_e = \min\nolimits_\delta\!\left(\frac{v_j}{N V_T},\; 50\right),
\end{equation}
realizzato con la \eqref{eq:smoothmin}. Con $N=2$ e $V_T = 26\,$mV il
clamp corrisponde a $v_j \approx 2{,}6\,$V, mai raggiunto in
funzionamento normale: la protezione non altera i risultati.

\paragraph{Valori nominali e scalatura del jacobiano.}
Gli stati di carica (ordine $10^{-7}$--$10^{-4}\,$C) convivono nella
stessa rete con tensioni e correnti di ordine $10^{2}$--$10^{3}$: senza
riscalatura il jacobiano di Newton copre oltre dieci ordini di
grandezza e pu\`o risultare mal condizionato, con errori del tipo
\emph{nonlinear solver failed to converge due to Linear Algebra error}.
Per questo tutte le variabili di carica dichiarano un valore
\ssc{nominal} esplicito, e lo stato della corrente di coda dell'IGBT
(sez.~\ref{sec:tail}) \`e tenuto direttamente in ampere anzich\'e come
carica in coulomb.

\section{Modello del diodo: \ssc{diode\_rr}}

\subsection{Modello a controllo di carica di Lauritzen--Ma}

Il modello \`e quello a due cariche concentrate di Lauritzen e Ma
\cite{lauritzen1991}. La carica in eccesso nella regione di drift \`e
rappresentata da due nodi: $q_E$, la carica al bordo della giunzione,
imposta istantaneamente dalla legge di giunzione, e $q_M$, la carica
immagazzinata al centro della regione. Tra i due nodi scorre una
corrente di diffusione proporzionale al gradiente di carica; $q_M$
decade per ricombinazione con la vita media $\tau$:
\begin{align}
i_{\mathrm{diff}} &= \frac{q_E - q_M}{T_M},
\label{eq:lm1}\\[2pt]
\frac{\ud q_M}{\ud t} &= \frac{q_E - q_M}{T_M} - \frac{q_M}{\tau},
\label{eq:lm2}\\[2pt]
q_E &= \tau\, I_s \left( e^{\,v_j/(N V_T)} - 1 \right).
\label{eq:lm3}
\end{align}
$T_M$ \`e il tempo di transito per diffusione attraverso la regione di
drift. La tensione ai morsetti e la corrente totale sono
\begin{equation}
v = v_j + R_s\, i,
\qquad
i = i_{\mathrm{diff}} + \frac{\ud q_j}{\ud t},
\end{equation}
con $q_j$ la carica di svuotamento (sez.~\ref{sec:depl}).

\subsection{Regime stazionario}

In continua la \eqref{eq:lm2} d\`a $q_M = \tau\, i$ e, sostituendo nella
\eqref{eq:lm1},
\begin{equation}
i = \frac{\tau}{\tau + T_M}\, I_s
\left( e^{\,v_j/(N V_T)} - 1 \right):
\label{eq:static}
\end{equation}
la caratteristica statica \`e una legge esponenziale con corrente di
saturazione efficace $I_s\,\tau/(\tau+T_M)$. Su questa curva, pi\`u la
caduta $R_s i$, si fittano $I_s$, $N$ e $R_s$ dalla caratteristica
diretta a caldo del datasheet.

\subsection{Reverse recovery ed estrazione di $\tau$ e $T_M$}
\label{sec:rr}

Durante la commutazione il circuito esterno impone una rampa di corrente
$i(t) = I_F - a\,t$ con $a = \ud i/\ud t$. Finch\'e la giunzione resta
polarizzata direttamente ($q_E > 0$, fase di storage) il diodo \`e a
bassa impedenza e la rampa prosegue anche sotto zero. Sostituendo
$q_E - q_M = T_M\, i$ nella \eqref{eq:lm2},
\begin{equation}
\frac{\ud q_M}{\ud t} = i - \frac{q_M}{\tau},
\end{equation}
che con $i(t) = I_F - a t$ e condizione iniziale $q_M(0) = \tau I_F$ ha
soluzione
\begin{equation}
q_M(t) = \tau\,(I_F - a t) + a \tau^2 \left(1 - e^{-t/\tau}\right).
\end{equation}
La fase di storage termina all'istante $t_s$ in cui $q_E = q_M + T_M i$
si annulla, cio\`e quando la giunzione inverte la polarizzazione; in quel
punto la corrente raggiunge il picco inverso $I_{rr} = a\,t_s - I_F$.
Imponendo $q_M(t_s) = -T_M\, i(t_s)$ si ottiene, per $t_s \gg \tau$,
\begin{equation}
\boxed{\;I_{rr} \simeq \frac{a\,\tau^2}{\tau + T_M}.\;}
\label{eq:irr}
\end{equation}
Dopo il picco $q_E \simeq 0$ e la \eqref{eq:lm2} diventa
$\dot q_M = -q_M \left( 1/T_M + 1/\tau \right)$: corrente e carica
decadono esponenzialmente (la coda del recovery) con costante di tempo
\begin{equation}
\boxed{\;\tau_{rr} = \frac{\tau\, T_M}{\tau + T_M}.\;}
\label{eq:taurr}
\end{equation}
La carica di recovery \`e la somma del triangolo fino al picco e della
coda:
\begin{equation}
Q_{rr} \simeq \frac{I_{rr}^2}{2a} + I_{rr}\,\tau_{rr}.
\label{eq:qrr}
\end{equation}

Invertendo le \eqref{eq:irr}--\eqref{eq:qrr} si ottiene la procedura di
estrazione in forma chiusa dai dati di targa ($I_{rr}$, $Q_{rr}$ a
$\ud i/\ud t = a$, alla corrente e temperatura di riferimento):
\begin{equation}
\boxed{\;
\tau_{rr} = \frac{Q_{rr}}{I_{rr}} - \frac{I_{rr}}{2a},
\qquad
\tau = \tau_{rr} + \frac{I_{rr}}{a},
\qquad
T_M = \frac{\tau\,\tau_{rr}\,a}{I_{rr}}.
\;}
\label{eq:extract}
\end{equation}
Verifica di consistenza: deve risultare $Q_{rr} > I_{rr}^2/(2a)$
(altrimenti $\tau_{rr}<0$); il rapporto $\tau/T_M$ controlla la softness
(coda lunga per $\tau \gg T_M$, recovery snappy per $\tau \lesssim T_M$).

\paragraph{Limite del modello.}
Nella \eqref{eq:irr} $I_{rr}$ cresce linearmente con $a$ ed \`e
indipendente da $I_F$ (per $t_s \gg \tau$); nel dispositivo reale la
crescita con $a$ \`e sublineare e c'\`e una dipendenza debole da $I_F$.
I parametri estratti valgono quindi \emph{attorno al punto di lavoro}
usato per l'estrazione: conviene tarare su $I_F$ e $\ud i/\ud t$
rappresentativi della commutazione che si vuole studiare.

\subsection{Capacit\`a di svuotamento}
\label{sec:depl}

Dopo il recovery la giunzione si polarizza inversamente e la dinamica
\`e dominata dalla capacit\`a di svuotamento, che insieme all'induttanza
di loop determina il ringing e il $\ud v/\ud t$ sul diodo. Si usa la
convenzione SPICE \cite{massobrio}:
\begin{equation}
C_j(v_j) = \frac{C_{j0}}{\left(1 - v_j/V_j\right)^{m}},
\qquad
q_j(v_j) = \frac{C_{j0} V_j}{1-m}
\left[ 1 - \left(1 - \frac{v_j}{V_j}\right)^{1-m} \right].
\end{equation}
La formula diverge per $v_j \to V_j$; come in SPICE, sopra
$v_j = F_C\, V_j$ (con $F_C = 0{,}5$) la si sostituisce con capacit\`a
costante $C_{FC} = C_{j0}/(1-F_C)^m$, cio\`e con l'estensione lineare
della carica. Nell'implementazione questo \`e realizzato senza
discontinuit\`a con il clamp smooth
$v_{jc} = \min_\delta(v_j,\, F_C V_j)$ e
\begin{equation}
q_j = \frac{C_{j0} V_j}{1-m}
\left[ 1 - \left(1 - \frac{v_{jc}}{V_j}\right)^{1-m} \right]
+ C_{FC}\,(v_j - v_{jc}),
\end{equation}
dove il secondo termine \`e nullo sotto il clamp ed \`e l'estensione
lineare sopra. La corrente capacitiva \`e $\ud q_j/\ud t$: la
formulazione in carica conserva la carica, a differenza di
$i = C(v)\,\ud v/\ud t$.

In diretta la capacit\`a di svuotamento \`e comunque mascherata dalla
capacit\`a di diffusione implicita nelle \eqref{eq:lm1}--\eqref{eq:lm3}.

\subsection{Parametri e procedura di taratura}

\begin{center}
\begin{tabular}{llll}
\hline
Simbolo & \ssc{.ssc} & Significato & Taratura \\
\hline
$I_s$   & \ssc{Is}  & corrente di saturazione & caratteristica statica \eqref{eq:static} \\
$N$     & \ssc{N}   & coefficiente di emissione & caratteristica statica \\
$V_T$   & \ssc{Vt}  & tensione termica $kT/q$ & $26$--$34\,$mV secondo $T_j$ \\
$\tau$  & \ssc{tau} & vita media dei portatori & \eqref{eq:extract} \\
$T_M$   & \ssc{TM}  & tempo di transito & \eqref{eq:extract} \\
$R_s$   & \ssc{Rs}  & resistenza serie & pendenza ad alta corrente \\
$C_{j0}$& \ssc{Cj0} & capacit\`a di svuotamento a $0\,$V & curva $C(V)$ \\
$V_j$   & \ssc{Vjp} & potenziale di giunzione & curva $C(V)$ \\
$m$     & \ssc{mg}  & coefficiente di grading & curva $C(V)$ \\
$F_C$   & \ssc{FC}  & clamp dello svuotamento & default $0{,}5$ \\
$\delta$& \ssc{dsm} & larghezza di smoothing & default $10\,$mV \\
\hline
\end{tabular}
\end{center}

Procedura consigliata: (1) $R_s$, $I_s$, $N$ dalla caratteristica
diretta a caldo; (2) $\tau$, $T_M$ dalle \eqref{eq:extract} con i dati
di recovery; (3) $C_{j0}$, $V_j$, $m$ dalla curva $C(V)$ del datasheet
oppure, a ritroso, dalla frequenza di ringing misurata; (4) verifica in
double pulse test confrontando $I_{rr}$ e $t_{rr}$ simulati con la
targa.

\section{Modello dell'IGBT: \ssc{igbt\_dyn}}

Il modello \`e comportamentale: canale MOS a legge quadratica, tre
capacit\`a interelettrodiche non lineari, resistenza interna di gate e
una corrente di coda a controllo di carica. Non \`e un modello fisico
alla Hefner \cite{hefner1994}, ma riproduce i meccanismi che determinano
i transitori: plateau di Miller, $\ud i/\ud t$ e $\ud v/\ud t$ imposti
dal circuito di gate, tail current allo spegnimento.

\subsection{Canale MOS}

Con $v_{ov} = v_{ge} - V_{th}$ (overdrive) e le parti positive
regolarizzate $v_{ov}^{+} = p_\delta(v_{ov})$,
$v_{ce}^{+} = p_\delta(v_{ce})$, si definisce
$v_x = \min_\delta\!\left(v_{ce}^{+},\, v_{ov}^{+}\right)$ e
\begin{equation}
i_{ch,0} = K_p \left( v_{ov}^{+}\, v_x - \frac{v_x^2}{2} \right)
\left( 1 + \lambda\, v_{ce}^{+} \right).
\label{eq:channel}
\end{equation}
Per $v_x = v_{ce}$ si ottiene la regione lineare
$K_p ( v_{ov} v_{ce} - v_{ce}^2/2 )$, per $v_x = v_{ov}$ la saturazione
$K_p v_{ov}^2 / 2$, con transizione smooth. $\lambda$ \`e la
conduttanza d'uscita (default $0$). Il canale \`e unidirezionale
($v_x \ge 0$): per $v_{ce} < 0$ conduce il \ssc{diode\_rr} esterno.

$V_{th}$ e $K_p$ si fittano sulla transcaratteristica a caldo,
$I_c = K_p (v_{ge}-V_{th})^2/2$; equivalentemente
$K_p = g_{fs}^2/(2 I_c)$ nel punto in cui il datasheet d\`a $g_{fs}$.
Una verifica utile del fit \`e la tensione di plateau di Miller,
\begin{equation}
v_{ge,\mathrm{pl}} = V_{th} + \sqrt{\frac{2 I_c}{K_p}},
\end{equation}
confrontabile con le forme d'onda di gate charge del datasheet.

\subsection{Capacit\`a non lineari}

Tutte le capacit\`a sono in formulazione di carica, $i = \ud Q/\ud t$.

\paragraph{$C_{ge}$ (costante).}
$Q_{ge} = C_{ge}\, v_{ge}$, con
$C_{ge} = C_{ies} - C_{res}$ valutata a $V_{ce}$ alto.

\paragraph{Capacit\`a di Miller $C_{gc}$.}
\`E la capacit\`a che governa plateau e $\ud v/\ud t$. Ha due regimi
fisici, distinti dal segno di $v_{gc}$:
\begin{itemize}
\item $v_{gc} > 0$ (dispositivo saturato, $v_{ce}$ basso): sotto
l'ossido di gate c'\`e accumulazione, la capacit\`a \`e quella
dell'ossido, costante e alta: $C_{gc} = C_{gc0}$;
\item $v_{gc} < 0$ (blocking): sotto l'ossido si forma una regione
svuotata in serie, la capacit\`a crolla con la tensione:
$C_{gc} = C_{gc0} / (1 - v_{gc}/V_{jc})^{m_c}$.
\end{itemize}
Separando $v_{gc}$ in parte positiva $v_g^{+} = p_\delta(v_{gc})$ e
negativa $v_g^{-} = v_{gc} - v_g^{+}$, la carica \`e
\begin{equation}
Q_{gc} = C_{gc0}\, v_g^{+}
- \frac{C_{gc0} V_{jc}}{1-m_c}
\left[ \left(1 - \frac{v_g^{-}}{V_{jc}}\right)^{1-m_c} - 1 \right].
\end{equation}
Per il fit: $C_{gc0} = C_{res}$ a $V_{ce} \approx 0$; $V_{jc}$ e $m_c$
si fittano sulla curva $C_{res}(V_{ce})$ del datasheet (per
$V_{ce} \gg v_{ge}$ vale $v_{gc} \approx -V_{ce}$). Le curve reali
scendono di 1--2 decadi tra $0$ e $25\,$V: servono tipicamente
$m_c \approx 0{,}6$--$0{,}9$ con $V_{jc} \approx 1\,$V.

\paragraph{Capacit\`a d'uscita $C_{ce}$.}
Svuotamento per $v_{ce} > 0$, costante $C_{ce0}$ sotto:
\begin{equation}
Q_{ce} = \frac{C_{ce0} V_{jo}}{1-m_o}
\left[ \left(1 + \frac{v_{ce}^{+}}{V_{jo}}\right)^{1-m_o} - 1 \right]
+ C_{ce0}\, (v_{ce} - v_{ce}^{+}).
\end{equation}
Si fitta su $C_{oes}(V_{ce}) - C_{res}(V_{ce})$. \`E la capacit\`a che,
con la $L_\sigma$ di loop e la $C_j$ del diodo, fissa la frequenza di
ringing dopo lo spegnimento.

\subsection{Corrente di coda}
\label{sec:tail}

Allo spegnimento di un IGBT il canale MOS si chiude in tempi imposti dal
gate, ma la carica dei portatori minoritari nella regione di drift
(iniettata dall'emettitore posteriore) deve ricombinare: una frazione
della corrente continua a scorrere e decade con la vita media dei
portatori. Il modello lo rappresenta con uno stato del primo ordine
sulla corrente di coda, tenuto direttamente in ampere per il
condizionamento numerico:
\begin{equation}
\tau_T\, \frac{\ud i_{\mathrm{tail}}}{\ud t}
= f_{\mathrm{tail}}\, i_{ch,0} - i_{\mathrm{tail}},
\qquad
i_C = (1 - f_{\mathrm{tail}})\, i_{ch,0} + i_{\mathrm{tail}}.
\label{eq:tail}
\end{equation}
In regime la \eqref{eq:tail} d\`a
$i_{\mathrm{tail}} = f_{\mathrm{tail}}\, i_{ch,0}$ e quindi
$i_C = i_{ch,0}$: la ripartizione non altera la statica. Allo
spegnimento $i_{ch,0}$ collassa in tempi di gate e la coda parte da
$f_{\mathrm{tail}}\, I_0$ decadendo con $\tau_T$. Valori tipici:
$f_{\mathrm{tail}} = 0{,}15$--$0{,}3$, $\tau_T = 1$--$3\,\mu$s, da
raffinare su una forma d'onda di $E_{off}$ del datasheet o su una
misura.

Limite: nella realt\`a la coda dipende anche da $v_{ce}$ (estrazione
della carica da parte della regione di carica spaziale); qui il
decadimento \`e puramente esponenziale. Per lo studio di overshoot e
snubber \`e adeguato, perch\'e la sovratensione \`e dominata dal
$\ud i/\ud t$ della caduta principale di corrente, non dalla coda.

\subsection{Rete di gate}

Il componente include la sola resistenza interna $R_g$ (dato di
datasheet); driver, $R_{g,on}$/$R_{g,off}$ esterne ed eventuale
induttanza di gate $L_g$ ($20$--$50\,$nH) vanno modellati fuori. Il
$\ud i/\ud t$ allo spegnimento emerge dalla scarica del gate attraverso
$R_{g,\mathrm{tot}}$ e la capacit\`a di Miller; durante il plateau
\begin{equation}
\frac{\ud v_{ce}}{\ud t} \approx \frac{i_g}{C_{gc}(v_{gc})},
\qquad
i_g = \frac{v_{ge,\mathrm{pl}} - V_{drv}}{R_{g,\mathrm{tot}}},
\end{equation}
il che permette di studiare direttamente il compromesso
$R_{g,off}$--sovratensione.

\subsection{Parametri}

\begin{center}
\begin{tabular}{llll}
\hline
Simbolo & \ssc{.ssc} & Significato & Taratura \\
\hline
$V_{th}$ & \ssc{Vth}  & tensione di soglia & transcaratteristica \\
$K_p$    & \ssc{Kp}   & parametro di transconduttanza & transcaratteristica, $g_{fs}$ \\
$\lambda$& \ssc{lam}  & conduttanza d'uscita & default $0$ \\
$C_{ge}$ & \ssc{Cge}  & capacit\`a gate--emettitore & $C_{ies}-C_{res}$ \\
$C_{gc0}$& \ssc{Cgc0} & Miller a $v_{gc}=0$ & $C_{res}$ a $V_{ce}\approx 0$ \\
$V_{jc}$, $m_c$ & \ssc{Vjc}, \ssc{mc} & svuotamento Miller & fit di $C_{res}(V_{ce})$ \\
$C_{ce0}$& \ssc{Cce0} & uscita a $v_{ce}=0$ & $C_{oes}-C_{res}$ a $V_{ce}\approx 0$ \\
$V_{jo}$, $m_o$ & \ssc{Vjo}, \ssc{mo} & svuotamento d'uscita & fit di $C_{oes}-C_{res}$ \\
$R_g$    & \ssc{Rg}   & resistenza interna di gate & datasheet \\
$f_{\mathrm{tail}}$ & \ssc{ftail} & frazione di coda & $E_{off}$ o misura \\
$\tau_T$ & \ssc{tauT} & costante di tempo della coda & $E_{off}$ o misura \\
$\delta$ & \ssc{dsm}  & larghezza di smoothing & default $0{,}1\,$V \\
\hline
\end{tabular}
\end{center}

\section{Uso in circuito e limiti}

Il banco di riferimento \`e il double pulse test: IGBT in basso come
dispositivo in prova, \ssc{diode\_rr} in alto, induttanza di carico in
parallelo al diodo, driver $\pm 15/{-}8\,$V con $R_{g,on}$ e $R_{g,off}$
separate. Il primo impulso serve anche da precondizionamento: porta la
corrente a regime e assesta $q_M$ del diodo prima della commutazione di
interesse.

La rete del DC link va modellata esternamente, dal condensatore verso il
modulo: banco bulk ($C_{dc}$ + ESR + ESL), busbar ($L_{bus}$ con
resistenza valutata alla frequenza di ringing, non in continua,
altrimenti l'oscillazione non si smorza realisticamente), e induttanza
interna del modulo $L_{sCE}$ in serie al collettore, tenuta separata
perch\'e nessuno snubber esterno pu\`o ridurla. Con
$L_\sigma = \mathrm{ESL} + L_{bus} + L_{sCE}$:
\begin{equation}
\Delta V_{ce} \approx L_\sigma \frac{\ud i}{\ud t},
\qquad
f_0 = \frac{1}{2\pi \sqrt{L_\sigma\, C_{eq}}},
\qquad
C_{eq} \approx C_{oes} + C_{j,\mathrm{diodo}};
\end{equation}
la frequenza di ringing \`e un ottimo cross check per tarare
$L_\sigma$ contro una misura.

Fenomeni riprodotti: plateau di Miller e gate charge; $\ud i/\ud t$ e
$\ud v/\ud t$ dipendenti da $R_g$; $I_{rr}$ e softness del diodo
dipendenti dal $\ud i/\ud t$ imposto; sovratensione
$L_\sigma\,\ud i/\ud t$ allo spegnimento; ringing e sua interazione con
snubber RC, RCD e condensatori di decoupling.

Limiti principali: nessun calcolo di perdite n\'e dipendenza dalla
temperatura (parametrizzare a caldo); coda non dipendente da $v_{ce}$;
modello di recovery valido attorno al punto di estrazione
(sez.~\ref{sec:rr}); forward recovery del diodo non riprodotto ($R_s$
costante), quindi il contributo di $V_{FR}$ alla sovratensione di
spegnimento dell'IGBT va eventualmente stimato a parte; capacit\`a
fittate localmente sulle curve del datasheet.

\begin{thebibliography}{9}

\bibitem{lauritzen1991}
P.~O. Lauritzen, C.~L. Ma,
\emph{A simple diode model with reverse recovery},
IEEE Transactions on Power Electronics, vol.~6, n.~2, pp.~188--191, 1991.

\bibitem{ma1993}
C.~L. Ma, P.~O. Lauritzen,
\emph{A simple power diode model with forward and reverse recovery},
IEEE Transactions on Power Electronics, vol.~8, n.~4, pp.~342--346, 1993.

\bibitem{hefner1994}
A.~R. Hefner, D.~M. Diebolt,
\emph{An experimentally verified IGBT model implemented in the Saber
circuit simulator},
IEEE Transactions on Power Electronics, vol.~9, n.~5, pp.~532--542, 1994.

\bibitem{massobrio}
G.~Massobrio, P.~Antognetti,
\emph{Semiconductor Device Modeling with SPICE}, 2a ed.,
McGraw--Hill, 1993.

\bibitem{baliga}
B.~J. Baliga,
\emph{Fundamentals of Power Semiconductor Devices},
Springer, 2008.

\end{thebibliography}

\end{document}
